% !TEX root = ../main.tex
\section{Formal Preliminaries}
\label{sec:preliminaries}

I assume familiarity with propositional modal logic; this section fixes notation
and records the facts I need.

Let $\mathcal{L}$ be a propositional language with the usual Boolean connectives
and a unary operator $X$, to be interpreted variously as $K$ for knowledge, $B$
for belief, or $\Out_M$ for model assertion. The relevant rules and schemata:

\begin{center}
\begin{tabular}{@{}ll@{}}
\toprule
RE (congruentiality) & from $\models \phi \leftrightarrow \psi$, infer $\models X\phi \leftrightarrow X\psi$ \\
RM (monotonicity rule) & from $\models \phi \rightarrow \psi$, infer $\models X\phi \rightarrow X\psi$ \\
N (necessitation) & from $\models \phi$, infer $\models X\phi$ \\
K (distribution) & $X(\phi \rightarrow \psi) \rightarrow (X\phi \rightarrow X\psi)$ \\
M (simplification) & $X(\phi \wedge \psi) \rightarrow (X\phi \wedge X\psi)$ \\
C (agglomeration) & $(X\phi \wedge X\psi) \rightarrow X(\phi \wedge \psi)$ \\
D (consistency) & $X\phi \rightarrow \neg X \neg \phi$ \\
T (factivity) & $X\phi \rightarrow \phi$ \\
4 (positive introspection) & $X\phi \rightarrow XX\phi$ \\
5 (negative introspection) & $\neg X\phi \rightarrow X \neg X\phi$ \\
\bottomrule
\end{tabular}
\end{center}

A modal logic is \emph{normal} iff it contains K and is closed under N and
uniform substitution; RE then follows. Following \citet{chellas1980}, a logic is
\emph{classical} iff closed under RE, \emph{monotonic} iff closed under RM, and
\emph{regular} iff monotonic and containing C. The hierarchy matters because RE
is the weakest of these conditions. A logic that fails RE is not merely
non-normal; it falls below the weakest rung of the classical hierarchy, and is
standardly called \emph{hyperintensional}.

Semantically, in Kripke semantics $X\phi$ is true at world $w$ iff $\phi$ is
true at every world accessible from $w$. Every Kripke frame, whatever its
accessibility conditions, validates RE, N, and K. Frame conditions---reflexivity,
transitivity, euclideanness, seriality---add axioms on top of the normal base;
they never subtract. Hence:

\begin{fact}
\label{fact:kripke}
If an operator fails RE, K, or N as a matter of its intended interpretation, it
has no Kripke semantics, for any accessibility relation whatever.
\end{fact}

Two weaker frameworks accommodate non-normal operators. In \emph{neighborhood}
(minimal) semantics \citep{scott1970,montague1970,chellas1980}, each world $w$ is
assigned a set $N(w)$ of sets of worlds, and $X\phi$ is true at $w$ iff the
truth set $\lVert \phi \rVert$ is a member of $N(w)$. Further axioms correspond
to closure conditions on $N(w)$: closure under supersets for RM, containing the
unit for N, and so on. In \emph{impossible-worlds} semantics
\citep{rantala1982,nolan1997,priest2005,berto2019}, the model contains, alongside
possible worlds, non-normal worlds at which truth is assigned to formulas
directly, unconstrained by logical closure.

The difference between the two will matter in
Section~\ref{sec:hyperintensionality}, so it is worth recording now. Bare
neighborhood semantics validates RE, and validates it in every frame: since
neighborhoods are sets of \emph{truth sets}, and logically equivalent formulas
have the same truth set, $X\phi \leftrightarrow X\psi$ follows automatically
whenever $\phi$ and $\psi$ are equivalent. Neighborhood semantics is thus the
semantics \emph{for} classical logics---the right tool for dropping K and N, and
the wrong tool, without modification, for dropping RE. Impossible worlds do
better on this score, because logically equivalent formulas can differ in truth
value at an impossible world.

\begin{fact}
\label{fact:neighborhood}
Every neighborhood frame validates RE. To model an operator that fails RE, the
framework must be hyperintensionalised---neighborhoods taken over formulations or
other finer-grained content surrogates rather than over sets of worlds---or
impossible worlds must be invoked.
\end{fact}

Finally, the received idealisations. Epistemic logic since \citet{hintikka1962}
standardly models knowledge as a normal factive operator, and belief as KD45.
I will not lean on the specific choice of axioms for knowledge, and should note
that the frequent shorthand ``knowledge is S5'' overstates the consensus: axiom
5 for knowledge, which delivers negative introspection, is widely rejected, and
Hintikka's own system was S4 \citep{lenzen1978}. Nothing in my argument requires
more than \emph{normality}, which is common to the whole family.

These idealisations are, notoriously, idealisations; the logical omniscience
problem is as old as the framework \citep{fagin1995}. But the tradition's
response has been to treat closure failures in humans as deviations from a
normal-operator competence---a point that bears weight in
Section~\ref{sec:human-parallel}. What the tradition has never countenanced as
knowledge or belief is an attitude whose constitutive profile is non-classical:
an attitude that is not even congruential in principle.

One omission is deliberate. I do not claim that $\Out_M$ fails C. Any
threshold-based account of assertion or belief fails agglomeration, for the
familiar lottery and preface reasons \citep{makinson1965}: contents may each
clear a threshold while their conjunction does not. That failure is generic to
thresholds, applies equally to threshold accounts of human belief, and is
therefore no evidence of anything distinctive about machines. I also do not rest
anything on N. There are infinitely many validities and finitely many outputs,
and the interesting question---whether models endorse feasibly checkable
tautologies under unfamiliar syntactic dress---is one on which I am not aware of
evidence clean enough to carry argumentative weight. Both principles are listed
above for completeness and then set aside.
